\section{Educational tracking in the Netherlands}\label{sec:Setting}

In the Netherlands, the transition to tracked secondary education takes place around the age of twelve when students leave the sixth and final grade of primary education. The secondary school system consists of three main tracks: (i) a university track (\textit{vwo}), which has a duration of six years and gives students access to university education, (ii) a college track (\textit{havo}), which has a duration of five years and prepares students for college education, and (iii) a vocational track (\textit{vmbo}), which has a duration of four years and serves as a preparation for vocational education. The vocational track consists of four sub-tracks, which can be ordered from more practice-oriented to more theory-oriented. 

The allocation of students to tracks is based on the recommendation provided by the sixth-grade primary school teacher. The recommendation process consists of two main steps. First, by March of the sixth grade the teacher provides an initial track recommendation. Then, in April or May, each student takes a standardized end-of-primary education test, which the teacher may use to consider an upward revision \citep{wpo2014}. Hence, the final track recommendation may be higher (but not lower) than the initial one. The teacher track recommendation is binding in the sense that, as a rule, secondary schools cannot place students in a higher track than indicated by the recommendation \citep{wvo2014}.

We use the initial, instead of the final, track recommendation as our measure for the teacher's subjective evaluation. This choice reflects that the initial recommendation is often the most influential, as many secondary schools begin allocating students to first-year classes as early as March, relying on these initial recommendations. Secondary schools may face challenges when changing their allocation by the time the final track recommendation is provided in May \citep{inspectorate2019}. Moreover, in practice, teachers upgrade only 1 in 10 students, meaning the initial and final track recommendations are identical for 9 out of 10 students \citep{deree2023quality}.

Teachers are expected to make holistic track recommendations by considering both objective standardized test results and their subjective judgments on factors such as motivation, attitude towards school, and classroom behavior. While there are no strict national rules describing how to combine these types of information, general guidelines from the \cite{ministry2022} suggest that standardized test results provide a ``good starting point''. A recent survey among 400 primary school teachers shows that 85\% of the respondents attach ``considerable'' to ``a lot of'' value to standardized test results when formulating their recommendations \citep{Duo2023}. 

\subsection{Standardized test scores}\label{sec:Testscores}

All Dutch primary school students, from grade one to six, are expected to take two standardized tests per year. The first test is administered in the middle of the school year around February and the second at the end of the school year in June. We will refer to these tests as the midterm and end term test, respectively. We use the fifth-grade end term test as the main measure of student ability to estimate the conditional SES gap in track recommendations. This test can be considered the last score that teachers observe before making the initial recommendation. Alternatively, one could use the end-of-primary school test score (the sixth-grade end term test) as the ability measure and the final track recommendation as the teacher's subjective evaluation. This alternative approach is problematic since the final track recommendation can only be adjusted upwards based on performance on the end-of-primary school test. Hence, this test is low (high) stakes for students who are (un)satisfied with the initial recommendation.

Primary schools have to select one of the standardized test systems that are approved by the Ministry of Education. Most schools opt for the test system provided by \textit{Cito}, which has a market share of roughly 80\% in the period we study. The biannual \textit{Cito} tests focus on two domains: mathematics and reading. The tests consist of approximately 70 to 100 items, include both multiple-choice and open-ended questions, and are scored by the teacher using standardized answer keys. The item scores are entered into the \textit{Cito} software and used to generate students' test scores through an item response theory (IRT) model. 

An IRT model specifies the probability that a student answers a test item correctly as a function of latent student ability and item parameters. The IRT model used by \textit{Cito} is the Birnbaum model \citep{Cito2017}, which includes two item parameters: difficulty and discrimination. The item difficulty parameter reflects the student ability level at which there is a 50\% probability of answering the item correctly. The discrimination parameter reflects how sharply the item distinguishes between students of different ability levels. These parameters may be estimated for a selection of items to form a so-called calibrated item bank, which can be used to create multiple tests of varying difficulty that aim to measure student ability on a common scale. This enables, for instance, observing the development of student ability from grade 1 to 6, one of the main goals of standardized testing in the Netherlands. 

\textit{Cito} constructs its calibrated item bank through national norming studies. First, sample tests are administered to representative student populations. Second, the item parameters are estimated: the discrimination parameters are held fixed while the difficulty parameters are estimated using conditional maximum likelihood. Both parameters are then iteratively adjusted to improve model fit. The resulting calibrated item bank enables the construction of the biannual tests for various grade levels. Although more difficult items are used in tests for higher grades, student ability can supposedly be measured on a common scale.

Latent student ability is estimated from these biannual tests as a ``fixed effect'' using weighted maximum likelihood \citep{sanders1993psychometrie}. In this approach, the two item parameters are held fixed at their calibrated values, and the likelihood function is weighted by the test information function. This function indicates how much information the test provides about student ability, with items of difficulty close to the student’s ability offering more information. \cite{warm1989weighted} shows that this procedure yields (nearly) unbiased estimates of student ability. The estimated ability scores obtained through this procedure are the test scores we use in our analysis. Unbiased, yet noisy, estimates of ability implies that the measurement error in the biannual \textit{Cito} test scores is assumed to be classical.